\section{Resultados I --- Descritivas do corpus}\label{cap:descritivas}

Este capítulo caracteriza o corpus de análise antes de passar à síntese
substantiva. Das 852 entradas que percorreram a extração estruturada, 816
constituem o conjunto efetivamente analisável --- estudos marcados como
elegíveis e com extração bem-sucedida. Excluem-se 35 registros inelegíveis (27
por tipo de documento inválido, E3, e 8 por qualidade insuficiente, E5) e um
único registro cuja resposta do modelo não pôde ser interpretada (\textit{parse
fail}). Salvo indicação em contrário, os percentuais deste e dos
próximos capítulos têm como base esse $N = 816$, e as proporções de cada
dimensão são calculadas apenas sobre os estudos que a classificaram (valores
não informados ficam fora do denominador).

\subsection{Volume e evolução temporal}

A Figura~\ref{fig:corpus_anos} mostra que o corpus se concentra fortemente no
período recente: 636 dos 816 estudos (77,9\%) foram publicados a partir de
2022, com pico em 2025 (252 estudos) e em 2024 (137); os 123 registros de 2026
correspondem a publicações em \textit{early access}. A produção anterior a 2018
é residual (o estudo mais antigo do corpus é de 2015). Esse perfil é, em si, um
primeiro achado: o interesse econômico pela relação entre inteligência
artificial e emprego acelera acentuadamente após a difusão dos modelos
generativos, refletindo tanto o aumento de submissões quanto a janela de busca
estendida a meados de 2026.

\begin{figure}[htbp]
\centering
\includegraphics[width=0.85\linewidth]{figures/corpus_anos.pdf}
\caption{Distribuição do corpus por ano de publicação.}
\label{fig:corpus_anos}
\end{figure}

\subsection{Distribuição por janela temporal}

Agrupando os estudos nas três janelas de interesse
(Figura~\ref{fig:corpus_janelas}), observam-se 15 estudos na janela da
automação (\janelaUm{}), 225 na janela de \textit{machine learning} e
\textit{deep learning} (\janelaDois{}) e 576 na janela da IA generativa
(\janelaTres{}). A janela e a posição pré/pós-ChatGPT são atribuídas de forma
determinística pelo ano de publicação (Seção~\ref{sec:periodo}): são
pós-ChatGPT os estudos publicados de 2023 em diante. Esta última janela
concentra 70,6\% do corpus, o que condiciona a
leitura comparativa dos capítulos seguintes: o peso da evidência recente é
grande, e a comparação pré/pós-ChatGPT deve ser interpretada à luz desse
desbalanceamento amostral.

\begin{figure}[htbp]
\centering
\includegraphics[width=0.7\linewidth]{figures/corpus_janelas.pdf}
\caption{Corpus por janela temporal.}
\label{fig:corpus_janelas}
\end{figure}

\subsection{Tipos de estudo}

Quanto ao tipo de evidência (Figura~\ref{fig:corpus_tipos}), predominam os
estudos de evidência macro/setorial (218) e as revisões e \textit{surveys}
(188), seguidos por trabalhos teóricos e de modelagem (148), estudos de
exposição ocupacional (131) e análises de firma ou de plataformas de
\textit{freelancers} (125); seis estudos foram rotulados pelo extrator como de
``indivíduo'', categoria não prevista no esquema e tratada, nas análises
multivariadas, junto à evidência de firma. A heterogeneidade de desenhos --- de exercícios de
exposição baseados em medidas como o índice de Felten \parencite{felten2021} a
estudos de plataformas de trabalho \textit{online} --- justifica a opção metodológica por uma síntese
narrativo-estruturada, e não meta-analítica.

\begin{figure}[htbp]
\centering
\includegraphics[width=0.85\linewidth]{figures/corpus_tipo_estudo.pdf}
\caption{Tipos de estudo no corpus.}
\label{fig:corpus_tipos}
\end{figure}

\subsection{Tecnologia de IA focada}

A Figura~\ref{fig:corpus_tecnologia} distribui o corpus pela tecnologia de IA em
foco. A maior parte dos estudos trata a IA de forma genérica (319) ou foca a
automação clássica (253) e a combinação de robôs e IA (160). A IA generativa e
os \textit{LLMs} aparecem explicitamente em 54 estudos, e as abordagens
preditivas de \textit{machine learning} em 27. Vale notar que a presença ainda
modesta da rotulagem ``IA generativa'' não contradiz a concentração temporal
recente: muitos estudos publicados após 2022 seguem analisando a automação e a
IA de modo amplo, e apenas uma fração nomeia diretamente os modelos generativos
como objeto. Essa constatação tem consequência analítica direta: entre os 576
estudos publicados de 2023 em diante, apenas 53 (9,2\%) têm a IA generativa como
objeto. ``Publicado depois do ChatGPT'' e ``sobre a tecnologia do ChatGPT'' são,
portanto, coisas distintas, e a Discussão (Capítulo~\ref{cap:discussao}) separa
as duas dimensões.

\begin{figure}[htbp]
\centering
\includegraphics[width=0.85\linewidth]{figures/corpus_tecnologia.pdf}
\caption{Distribuição do corpus por tecnologia de IA focada.}
\label{fig:corpus_tecnologia}
\end{figure}

\subsection{Atributos estruturais}

A Tabela~\ref{tab:descritivas_estrutural} resume atributos estruturais do
corpus. A maior parte são artigos de periódico (77,3\%) e 86,2\% são revisados
por pares, o que sustenta a qualidade editorial mínima do conjunto. Quanto à
estratégia empírica, predomina a análise descritiva (44,5\%), seguida de
regressões por \gls{mqo} (14,8\%), modelos teóricos (9,8\%), modelos estruturais
(9,7\%) e variáveis instrumentais (9,1\%); desenhos quasi-experimentais mais
exigentes em identificação causal (diferenças-em-diferenças, \gls{rdd}, estudo de
evento) somam menos de 5\%. Geograficamente, quase metade dos estudos é
multipaís (46,9\%), com a China (14,2\%) e os Estados Unidos à frente entre os
países isolados. A presença chinesa é um fenômeno recente: os estudos com foco na
China passam de 3,3\% do corpus pré-ChatGPT para 18,8\% do pós, de modo que a
expansão da literatura é também um deslocamento geográfico de quem a produz e de
onde ela olha. Nenhum \textit{working paper} integra o corpus, pois as bases
consultadas indexam periódicos e capítulos; a implicação para a cobertura é
discutida na Seção~\ref{sec:limitacoes}. Esse predomínio do descritivo sobre o identificado deve ser
mantido em vista na leitura dos achados de magnitude
(Seção~\ref{cap:comparacao}) e reaparece na análise de robustez por qualidade.

\begin{table}[htbp]
\centering
\caption{Atributos estruturais do corpus de análise.}
\label{tab:descritivas_estrutural}
\input{tables/descritivas_corpus.tex}
\end{table}
